\documentclass[10pt,twocolumn]{article}

\usepackage[margin=0.75in,top=0.7in,bottom=0.7in]{geometry}
\usepackage{amsmath,amssymb}
\usepackage{booktabs}
\usepackage{hyperref}
\usepackage{microtype}
\usepackage{parskip}
\usepackage{titlesec}
\usepackage{cite}

\titlespacing*{\section}{0pt}{6pt}{3pt}
\titlespacing*{\subsection}{0pt}{4pt}{2pt}
\setlength{\columnsep}{0.2in}
\setlength{\parskip}{3pt}

\title{\vspace{-1.5em}
\textbf{Physics-Guided Invariant Learning for\\LPBF Defect Segmentation}
\vspace{-0.5em}}

\author{
Ziming Zhao\\
\small Mechanical Engineering, Northwestern University\\
\small \texttt{zimingzhao2025@u.northwestern.edu}
}
\date{\vspace{-1em}}

\begin{document}
\maketitle
\thispagestyle{empty}

\begin{abstract}
\noindent
Three Scientific Machine Learning mechanisms are applied to LPBF
powder bed defect segmentation, each embedding the governing
Volumetric Energy Density equation $\text{VED}=P/(v\cdot h\cdot t)$
in a different way. NOTEARS causal structure learning recovers the
process graph $\{P,v,h,t\}\to\text{VED}\to\text{defect}$ from data.
Invariant Risk Minimisation (IRM) stratifies UNet training by VED
regime, yielding interpretable failure modes rather than a single
aggregate metric. A Physics-Informed Neural Network (PINN) encodes
the defect-nucleation ODE as a residual loss term, and a
physics-constrained Neural ODE models layer-wise defect evolution.
On the PB dataset (2,638 images, EOS M290, 316L~SS), the baseline
UNet achieves Dice\,=\,0.787 and ROC~AUC\,=\,0.9955.
The PINN ablation shows that a physics penalty weight of
$\lambda=1$ reduces MSE against the ODE ground truth relative to
a data-only baseline.
\end{abstract}

\section{Introduction}

In Laser Powder Bed Fusion, the dominant process variable is
Volumetric Energy Density:
\begin{equation}
  \text{VED} = \frac{P}{v \cdot h \cdot t},
  \label{eq:ved}
\end{equation}
where $P$ is laser power, $v$ scan speed, $h$ hatch spacing, and
$t$ layer thickness. Below a critical VED threshold
($\approx 30$\,J/mm$^3$ for 316L\,SS), powder fails to fully melt
and lack-of-fusion (LoF) pores form~\cite{debroy2018}.

Standard segmentation pipelines treat defect detection as a purely
data-driven problem. The physics is never consulted. The question
here is whether building Eq.~\eqref{eq:ved} into the learning
pipeline changes anything meaningful, and if so, how.

Three SciML mechanisms are tested. NOTEARS~\cite{zheng2018} is used
to recover the causal graph over process variables from data, without
being told the VED formula. IRM~\cite{arjovsky2019} is used to train
a UNet with a penalty that forces invariance across VED-defined
environments. A PINN~\cite{raissi2019} encodes the nucleation ODE as
a soft constraint evaluated at collocation points, requiring no
additional image labels. Each mechanism embeds the physics
differently, and the payoffs are compared.

\section{Method}

\subsection{UNet Baseline}

A standard four-level UNet~\cite{ronneberger2015} is used, with
channel depths $64\to128\to256\to512$ and a 1024-channel bottleneck.
Ten one-factor-at-a-time experiments vary learning rate
($10^{-4}$, $5\times10^{-4}$, $10^{-3}$), loss function (BCE, Dice,
BCEDice, Focal~\cite{lin2017}), optimiser (Adam, SGD, RMSprop), and
threshold $\tau\in\{0.3,0.5,0.7\}$. Batch size is 4, with 50 epochs
maximum and patience 7. The combined BCEDice loss is
\begin{equation}
  \mathcal{L} = \mathcal{L}_{\text{BCE}}
  + 1 - \frac{2\sum_i y_i\hat{p}_i + \varepsilon}
             {\sum_i y_i + \sum_i \hat{p}_i + \varepsilon}.
\end{equation}

\subsection{NOTEARS and IRM}

NOTEARS fits a weighted adjacency $W$ over
$(P, v, t, h, \text{VED}, \text{defect\_area}, \text{regime})$:
\begin{equation}
  \min_{W}\;\tfrac{1}{2n}\|X-XW\|_F^2 + \lambda_1\|W\|_1
  \quad \text{s.t.}\; \operatorname{tr}(e^{W\circ W}) - d = 0,
\end{equation}
where the acyclicity constraint rules out directed cycles.
After physics-consistency filtering (edges that violate the sign of
$\partial\text{VED}/\partial P$, etc., are removed), the recovered
structure is checked against Eq.~\eqref{eq:ved}.

VED values split the dataset into two IRM environments:
$e_1$ (stable, $30\leq\text{VED}\leq80$\,J/mm$^3$) and
$e_2$ (low, VED\,$<30$). Training minimises
\begin{equation}
  \min_{\phi}\;\sum_e \mathcal{L}^e(\phi)
  + \lambda_{\text{IRM}}
  \bigl\|\nabla_{w|w=1}\mathcal{L}^e(\phi,w)\bigr\|^2,
\end{equation}
with $\lambda_{\text{IRM}}$ annealed from $0$ to $10^3$.

\subsection{PINN and Neural ODE}

The defect-nucleation ODE is taken as
\begin{equation}
  \frac{df}{d\,\text{VED}} = -\alpha(f - f_{\text{eq}}(\text{VED})),
  \label{eq:ode}
\end{equation}
where $f_{\text{eq}} = f_{\max}\exp(-\beta(\text{VED}-\text{VED}_c)^2)$.
A four-layer Tanh MLP maps VED to defect fraction and is trained with
\begin{equation}
  \mathcal{L} = \mathcal{L}_{\text{data}}
  + \lambda_{\text{phys}}\,\mathcal{L}_{\text{ODE}},
\end{equation}
where $\mathcal{L}_{\text{ODE}}$ is the mean-squared residual of
Eq.~\eqref{eq:ode} at collocation points. An ablation over
$\lambda_{\text{phys}}\in\{0,0.1,1,10\}$ isolates the physics term.

The Neural ODE extends this to layer-wise dynamics,
\begin{equation}
  \frac{df}{dl} = -\gamma\,\text{VED}(l)\cdot f + s(\text{VED}(l)),
\end{equation}
comparing a pure-MLP vector field against one where the healing-rate
structure is fixed by the network architecture.

\section{Results}

\subsection{Hyperparameter Sweep}

\begin{table}[h]
\centering
\caption{Sweep results, four representative experiments.}
\label{tab:sweep}
\small
\setlength{\tabcolsep}{4pt}
\begin{tabular}{llcc}
\toprule
Loss & lr & Dice & ROC AUC \\
\midrule
BCEDice & $5\times10^{-4}$ & \textbf{0.787} & \textbf{0.9955} \\
Dice    & $1\times10^{-4}$ & 0.790          & 0.9941 \\
BCE     & $1\times10^{-4}$ & 0.766          & 0.9938 \\
Focal   & $1\times10^{-4}$ & 0.757          & 0.9921 \\
\bottomrule
\end{tabular}
\end{table}

Overlap-aware losses (BCEDice, Dice) outperform pixel-wise losses
under the roughly 1\% foreground fraction of the dataset. Pure Dice
gives marginally higher Dice, but BCEDice is more reliable in
practice because the BCE term penalises false positives on
non-defective powder texture. At image level, the best configuration
detects all 301 defective frames with zero misses and 4 false alarms.
Threshold variation from $\tau=0.3$ to $0.7$ shifts Dice by less
than 0.003, which confirms well-calibrated output probabilities.

\subsection{NOTEARS and IRM}

The recovered edges are: $t\to\text{VED}$ ($+1.00$),
$h\to\text{VED}$ ($+1.00$), $v\to\text{VED}$ ($-0.84$),
$P\to v$ ($-0.91$), $\text{VED}\to\text{defect\_area}$ ($+0.33$).
All signs match the analytic partial derivatives of
Eq.~\eqref{eq:ved}. The $P\to v$ edge reflects the common
operator practice of co-varying power and speed to hold a target VED.

IRM gives Dice\,=\,0.633 on stable-VED images and 0.525 on low-VED
images. The gap has a physical explanation: at stable VED the melt
pool re-solidifies smoothly and defects appear as bounded dark
regions, while at low VED the sintered powder retains inter-particle
voids at the same spatial scale as LoF pores.
The optimal threshold shifts from $\tau^*=0.85$ (stable) to
$0.90$ (low VED), since the model assigns moderate confidence scores
to ambiguous powder texture. This shift is not detectable without
regime-stratified evaluation.

\subsection{PINN}

\begin{table}[h]
\centering
\caption{PINN ablation: MSE against ODE ground truth.}
\label{tab:pinn}
\small
\begin{tabular}{cc}
\toprule
$\lambda_{\text{phys}}$ & MSE (vs.\ ODE) \\
\midrule
0.0  & highest \\
0.1  & reduced \\
1.0  & \textbf{lowest} \\
10.0 & slightly higher \\
\bottomrule
\end{tabular}
\end{table}

At $\lambda=0$ the network fits the noisy training observations but
oscillates outside the training range, violating the ODE.
At $\lambda=1$ the physics residual is small and the predicted
nucleation curve matches the ODE solution across the full VED range,
including the low-VED region where training data is sparse.
At $\lambda=10$ the physics constraint dominates and the network
underfits near the LoF boundary.
The physics-constrained Neural ODE converges faster than the
pure-MLP version and generalises to initial conditions outside
the training set, a direct result of the healing-rate structure
being fixed by the architecture.

\section{Discussion}

The three mechanisms test different ways of using the same piece of
physics. NOTEARS checks whether the governing equation is
recoverable from data at all, and the answer here is yes, with edge
signs consistent with Eq.~\eqref{eq:ved}. IRM uses the recovered
structure to define training environments and produces
regime-specific failure modes that trace back to process conditions.
The PINN shows that even a one-dimensional physics constraint,
evaluated at unlabelled collocation points, regularises the network
toward the correct nucleation curve.

The IRM Dice scores are lower than the baseline. This is expected:
the invariance penalty deliberately prevents the model from
exploiting VED-regime-specific texture patterns, so raw Dice drops.
The payoff is not metric gain but interpretability -- the model
fails in predictable, physically consistent ways across regimes.

Remaining limitations include the absence of high-VED keyhole data,
the use of synthetic process-variable data in the NOTEARS stage
(real per-image parameter measurements would strengthen the causal
analysis), and the gap between the 1-D PINN model and actual
image-level defect prediction. Closing these gaps is left for future
work.

\section*{References}
\begingroup
\small
\renewcommand{\section}[2]{}
\begin{thebibliography}{9}

\bibitem{debroy2018}
T.~DebRoy et al.,
``Additive manufacturing of metallic components,''
\textit{Prog.\ Mater.\ Sci.}, vol.~92, pp.~112--224, 2018.

\bibitem{ronneberger2015}
O.~Ronneberger, P.~Fischer, T.~Brox,
``U-Net,''
\textit{MICCAI}, pp.~234--241, 2015.

\bibitem{zheng2018}
X.~Zheng et al.,
``DAGs with NO TEARS,''
\textit{NeurIPS}, pp.~9472--9483, 2018.

\bibitem{arjovsky2019}
M.~Arjovsky et al.,
``Invariant risk minimization,''
\textit{arXiv:1907.02893}, 2019.

\bibitem{lin2017}
T.-Y.~Lin et al.,
``Focal loss for dense object detection,''
\textit{ICCV}, pp.~2980--2988, 2017.

\bibitem{raissi2019}
M.~Raissi, P.~Perdikaris, G.~Karniadakis,
``Physics-informed neural networks,''
\textit{J.\ Comput.\ Phys.}, vol.~378, pp.~686--707, 2019.

\bibitem{chen2018}
R.~T.~Q.~Chen et al.,
``Neural ordinary differential equations,''
\textit{NeurIPS}, 2018.

\end{thebibliography}
\endgroup

\end{document}
